\documentclass{article}
\usepackage{amsmath,amssymb}
\usepackage{bm}
\usepackage{geometry}
\usepackage{ctex}
\geometry{a4paper,margin=1.8cm}
\begin{document}

% 第1题
\noindent 1. 用Schmidt正交化方法求矩阵
\[
A=\begin{bmatrix}
1 & 2 & 2 \\
2 & 1 & 2 \\
1 & 2 & 1
\end{bmatrix}
\]
的QR分解。

\noindent \textbf{答：}
记 $A=[\bm{\alpha}_1,\bm{\alpha}_2,\bm{\alpha}_3]$，其中
\[
\bm{\alpha}_1=\begin{bmatrix}1\\2\\1\end{bmatrix},\quad
\bm{\alpha}_2=\begin{bmatrix}2\\1\\2\end{bmatrix},\quad
\bm{\alpha}_3=\begin{bmatrix}2\\2\\1\end{bmatrix}.
\]

\textbf{第一步：Schmidt正交化。}
\begin{align*}
\bm{\beta}_1 &= \bm{\alpha}_1 = \begin{bmatrix}1\\2\\1\end{bmatrix},\\[4pt]
\bm{\beta}_2 &= \bm{\alpha}_2 - \frac{\langle\bm{\alpha}_2,\bm{\beta}_1\rangle}{\langle\bm{\beta}_1,\bm{\beta}_1\rangle}\bm{\beta}_1
= \begin{bmatrix}2\\1\\2\end{bmatrix} - \frac{2\cdot1+1\cdot2+2\cdot1}{1^2+2^2+1^2}\begin{bmatrix}1\\2\\1\end{bmatrix}
= \begin{bmatrix}2\\1\\2\end{bmatrix} - \frac{6}{6}\begin{bmatrix}1\\2\\1\end{bmatrix}
= \begin{bmatrix}1\\-1\\1\end{bmatrix},\\[4pt]
\bm{\beta}_3 &= \bm{\alpha}_3 - \frac{\langle\bm{\alpha}_3,\bm{\beta}_1\rangle}{\langle\bm{\beta}_1,\bm{\beta}_1\rangle}\bm{\beta}_1
- \frac{\langle\bm{\alpha}_3,\bm{\beta}_2\rangle}{\langle\bm{\beta}_2,\bm{\beta}_2\rangle}\bm{\beta}_2\\
&= \begin{bmatrix}2\\2\\1\end{bmatrix} - \frac{2\cdot1+2\cdot2+1\cdot1}{6}\begin{bmatrix}1\\2\\1\end{bmatrix}
- \frac{2\cdot1+2\cdot(-1)+1\cdot1}{3}\begin{bmatrix}1\\-1\\1\end{bmatrix}\\
&= \begin{bmatrix}2\\2\\1\end{bmatrix} - \frac{7}{6}\begin{bmatrix}1\\2\\1\end{bmatrix} - \frac{1}{3}\begin{bmatrix}1\\-1\\1\end{bmatrix}
= \begin{bmatrix}\frac12\\[2pt]0\\[2pt]-\frac12\end{bmatrix}.
\end{align*}

\textbf{第二步：单位化。}
\begin{align*}
\bm{q}_1 &= \frac{\bm{\beta}_1}{\|\bm{\beta}_1\|} = \frac{1}{\sqrt{6}}\begin{bmatrix}1\\2\\1\end{bmatrix},
\quad \|\bm{\beta}_1\|=\sqrt{6},\\[4pt]
\bm{q}_2 &= \frac{\bm{\beta}_2}{\|\bm{\beta}_2\|} = \frac{1}{\sqrt{3}}\begin{bmatrix}1\\-1\\1\end{bmatrix},
\quad \|\bm{\beta}_2\|=\sqrt{3},\\[4pt]
\bm{q}_3 &= \frac{\bm{\beta}_3}{\|\bm{\beta}_3\|} = \sqrt{2}\begin{bmatrix}\frac12\\[2pt]0\\[2pt]-\frac12\end{bmatrix}
= \frac{1}{\sqrt{2}}\begin{bmatrix}1\\0\\-1\end{bmatrix},\quad \|\bm{\beta}_3\|=\frac{1}{\sqrt{2}}.
\end{align*}

\textbf{第三步：写出 $Q$ 和 $R$。}
\[
Q = \begin{bmatrix}
\frac{1}{\sqrt{6}} & \frac{1}{\sqrt{3}} & \frac{1}{\sqrt{2}}\\[4pt]
\frac{2}{\sqrt{6}} & -\frac{1}{\sqrt{3}} & 0\\[4pt]
\frac{1}{\sqrt{6}} & \frac{1}{\sqrt{3}} & -\frac{1}{\sqrt{2}}
\end{bmatrix},\qquad
R = Q^{\mathsf T}A = \begin{bmatrix}
\sqrt{6} & \sqrt{6} & \frac{7}{\sqrt{6}}\\[4pt]
0 & \sqrt{3} & \frac{1}{\sqrt{3}}\\[4pt]
0 & 0 & \frac{1}{\sqrt{2}}
\end{bmatrix}.
\]
\newpage

% 第2题
\noindent 2. 用Schmidt正交化方法求矩阵
\[
A=\begin{bmatrix}
0 & 1 & 1 \\
1 & 1 & 0 \\
1 & 0 & 1
\end{bmatrix}
\]
的QR分解。

\noindent \textbf{答：}
记 $A=[\bm{\alpha}_1,\bm{\alpha}_2,\bm{\alpha}_3]$，其中
\[
\bm{\alpha}_1=\begin{bmatrix}0\\1\\1\end{bmatrix},\quad
\bm{\alpha}_2=\begin{bmatrix}1\\1\\0\end{bmatrix},\quad
\bm{\alpha}_3=\begin{bmatrix}1\\0\\1\end{bmatrix}.
\]

\textbf{第一步：Schmidt正交化。}
\begin{align*}
\bm{\beta}_1 &= \bm{\alpha}_1 = \begin{bmatrix}0\\1\\1\end{bmatrix},\quad \|\bm{\beta}_1\|=\sqrt{2},\\[4pt]
\bm{\beta}_2 &= \bm{\alpha}_2 - \frac{\langle\bm{\alpha}_2,\bm{\beta}_1\rangle}{\langle\bm{\beta}_1,\bm{\beta}_1\rangle}\bm{\beta}_1
= \begin{bmatrix}1\\1\\0\end{bmatrix} - \frac{1}{2}\begin{bmatrix}0\\1\\1\end{bmatrix}
= \begin{bmatrix}1\\[2pt]\frac12\\[2pt]-\frac12\end{bmatrix},\quad
\|\bm{\beta}_2\|=\sqrt{1+\frac14+\frac14}=\frac{\sqrt{6}}{2},\\[4pt]
\bm{\beta}_3 &= \bm{\alpha}_3 - \frac{\langle\bm{\alpha}_3,\bm{\beta}_1\rangle}{\langle\bm{\beta}_1,\bm{\beta}_1\rangle}\bm{\beta}_1
- \frac{\langle\bm{\alpha}_3,\bm{\beta}_2\rangle}{\langle\bm{\beta}_2,\bm{\beta}_2\rangle}\bm{\beta}_2\\
&= \begin{bmatrix}1\\0\\1\end{bmatrix} - \frac{1}{2}\begin{bmatrix}0\\1\\1\end{bmatrix}
- \frac{\frac12}{\frac32}\begin{bmatrix}1\\[2pt]\frac12\\[2pt]-\frac12\end{bmatrix}
= \begin{bmatrix}1\\0\\1\end{bmatrix} - \begin{bmatrix}0\\[2pt]\frac12\\[2pt]\frac12\end{bmatrix}
- \begin{bmatrix}\frac13\\[2pt]\frac16\\[2pt]-\frac16\end{bmatrix}
= \begin{bmatrix}\frac23\\[2pt]-\frac23\\[2pt]\frac23\end{bmatrix},\quad
\|\bm{\beta}_3\|=\frac{2}{\sqrt{3}}.
\end{align*}

\textbf{第二步：单位化。}
\[
\bm{q}_1 = \frac{1}{\sqrt{2}}\begin{bmatrix}0\\1\\1\end{bmatrix},\qquad
\bm{q}_2 = \frac{1}{\sqrt{6}}\begin{bmatrix}2\\1\\-1\end{bmatrix},\qquad
\bm{q}_3 = \frac{1}{\sqrt{3}}\begin{bmatrix}1\\-1\\1\end{bmatrix}.
\]

\textbf{第三步：写出 $Q$ 和 $R$。}
\[
Q = \begin{bmatrix}
0 & \frac{2}{\sqrt{6}} & \frac{1}{\sqrt{3}}\\[4pt]
\frac{1}{\sqrt{2}} & \frac{1}{\sqrt{6}} & -\frac{1}{\sqrt{3}}\\[4pt]
\frac{1}{\sqrt{2}} & -\frac{1}{\sqrt{6}} & \frac{1}{\sqrt{3}}
\end{bmatrix},\qquad
R = \begin{bmatrix}
\sqrt{2} & \frac{1}{\sqrt{2}} & \frac{1}{\sqrt{2}}\\[4pt]
0 & \frac{\sqrt{6}}{2} & \frac{1}{\sqrt{6}}\\[4pt]
0 & 0 & \frac{2}{\sqrt{3}}
\end{bmatrix}.
\]
\newpage

% 第3题
\noindent 3. 用Givens变换求矩阵
\[
A=\begin{bmatrix}
0 & 1 & 1 \\
1 & 1 & 0 \\
1 & 0 & 1
\end{bmatrix}
\]
的QR分解。

\noindent \textbf{答：}
Givens变换通过一系列平面旋转变换逐列消去对角线下方元素。

\textbf{第一步：处理第一列。}
先消去 $a_{21}=1$：构造 $G(1,2,\theta_1)$，其中
$c_1=\frac{a_{11}}{\sqrt{a_{11}^2+a_{21}^2}}=0,\; s_1=\frac{a_{21}}{\sqrt{a_{11}^2+a_{21}^2}}=1$。
\[
G_1=\begin{bmatrix}
0 & 1 & 0\\-1 & 0 & 0\\0 & 0 & 1
\end{bmatrix},\quad
A^{(1)}=G_1A=\begin{bmatrix}
1 & 1 & 0\\0 & -1 & -1\\1 & 0 & 1
\end{bmatrix}.
\]

再消去 $a_{31}^{(1)}=1$：构造 $G(1,3,\theta_2)$，其中
$r=\sqrt{1^2+1^2}=\sqrt{2},\;c_2=\frac{1}{\sqrt{2}},\;s_2=\frac{1}{\sqrt{2}}$。
\[
G_2=\begin{bmatrix}
\frac{1}{\sqrt{2}} & 0 & \frac{1}{\sqrt{2}}\\[4pt]
0 & 1 & 0\\[4pt]
-\frac{1}{\sqrt{2}} & 0 & \frac{1}{\sqrt{2}}
\end{bmatrix},\quad
A^{(2)}=G_2A^{(1)}=\begin{bmatrix}
\sqrt{2} & \frac{1}{\sqrt{2}} & \frac{1}{\sqrt{2}}\\[4pt]
0 & -1 & -1\\[4pt]
0 & -\frac{1}{\sqrt{2}} & \frac{1}{\sqrt{2}}
\end{bmatrix}.
\]

\textbf{第二步：处理第二列。}
消去 $a_{32}^{(2)}=-\frac{1}{\sqrt{2}}$：构造 $G(2,3,\theta_3)$，其中
$r=\sqrt{(-1)^2+(-\frac{1}{\sqrt{2}})^2}=\frac{\sqrt{6}}{2},\;
c_3=\frac{-1}{r}=-\frac{2}{\sqrt{6}},\;s_3=\frac{-1/\sqrt{2}}{r}=-\frac{1}{\sqrt{3}}$。
\[
G_3=\begin{bmatrix}
1 & 0 & 0\\[4pt]
0 & -\frac{2}{\sqrt{6}} & -\frac{1}{\sqrt{3}}\\[4pt]
0 & \frac{1}{\sqrt{3}} & -\frac{2}{\sqrt{6}}
\end{bmatrix},\quad
R=A^{(3)}=G_3A^{(2)}=\begin{bmatrix}
\sqrt{2} & \frac{1}{\sqrt{2}} & \frac{1}{\sqrt{2}}\\[4pt]
0 & \frac{3}{\sqrt{6}} & \frac{1}{\sqrt{6}}\\[4pt]
0 & 0 & -\frac{2}{\sqrt{3}}
\end{bmatrix}.
\]

\textbf{第三步：求 $Q$。}
由 $Q^{\mathsf T}=G_3G_2G_1$，得 $Q=G_1^{\mathsf T}G_2^{\mathsf T}G_3^{\mathsf T}$：
\[
Q = \begin{bmatrix}
0 & -1 & 0\\1 & 0 & 0\\0 & 0 & 1
\end{bmatrix}
\begin{bmatrix}
\frac{1}{\sqrt{2}} & 0 & -\frac{1}{\sqrt{2}}\\[4pt]
0 & 1 & 0\\[4pt]
\frac{1}{\sqrt{2}} & 0 & \frac{1}{\sqrt{2}}
\end{bmatrix}
\begin{bmatrix}
1 & 0 & 0\\[4pt]
0 & -\frac{2}{\sqrt{6}} & \frac{1}{\sqrt{3}}\\[4pt]
0 & -\frac{1}{\sqrt{3}} & -\frac{2}{\sqrt{6}}
\end{bmatrix}
= \begin{bmatrix}
0 & \frac{2}{\sqrt{6}} & -\frac{1}{\sqrt{3}}\\[4pt]
\frac{1}{\sqrt{2}} & \frac{1}{\sqrt{6}} & \frac{1}{\sqrt{3}}\\[4pt]
\frac{1}{\sqrt{2}} & -\frac{1}{\sqrt{6}} & -\frac{1}{\sqrt{3}}
\end{bmatrix}.
\]
\newpage

% 第4题
\noindent 4. 用Givens变换求矩阵
\[
A=\begin{bmatrix}
2 & 2 & 1 \\
0 & 2 & 2 \\
2 & 1 & 2
\end{bmatrix}
\]
的QR分解。

\noindent \textbf{答：}

\textbf{第一步：处理第一列。}
$a_{21}=0$ 已为零，跳过。消去 $a_{31}=2$：构造 $G(1,3,\theta_1)$，
$r=\sqrt{2^2+2^2}=2\sqrt{2},\;c_1=s_1=\frac{1}{\sqrt{2}}$。
\[
G_1=\begin{bmatrix}
\frac{1}{\sqrt{2}} & 0 & \frac{1}{\sqrt{2}}\\[4pt]
0 & 1 & 0\\[4pt]
-\frac{1}{\sqrt{2}} & 0 & \frac{1}{\sqrt{2}}
\end{bmatrix},\quad
A^{(1)}=G_1A=\begin{bmatrix}
2\sqrt{2} & \frac{3}{\sqrt{2}} & \frac{3}{\sqrt{2}}\\[4pt]
0 & 2 & 2\\[4pt]
0 & -\frac{1}{\sqrt{2}} & \frac{1}{\sqrt{2}}
\end{bmatrix}.
\]

\textbf{第二步：处理第二列。}
消去 $a_{32}^{(1)}=-\frac{1}{\sqrt{2}}$：构造 $G(2,3,\theta_2)$，
$r=\sqrt{2^2+(-\frac{1}{\sqrt{2}})^2}=\frac{3}{\sqrt{2}},\;
c_2=\frac{2}{r}=\frac{2\sqrt{2}}{3},\;s_2=\frac{-1/\sqrt{2}}{r}=-\frac{1}{3}$。
\[
G_2=\begin{bmatrix}
1 & 0 & 0\\[4pt]
0 & \frac{2\sqrt{2}}{3} & -\frac{1}{3}\\[4pt]
0 & \frac{1}{3} & \frac{2\sqrt{2}}{3}
\end{bmatrix},\quad
R=A^{(2)}=G_2A^{(1)}=\begin{bmatrix}
2\sqrt{2} & \frac{3}{\sqrt{2}} & \frac{3}{\sqrt{2}}\\[4pt]
0 & \frac{3\sqrt{2}}{2} & \frac{7\sqrt{2}}{6}\\[4pt]
0 & 0 & \frac{4}{3}
\end{bmatrix}.
\]

\textbf{第三步：求 $Q$。}
$Q=G_1^{\mathsf T}G_2^{\mathsf T}$：
\[
Q = \begin{bmatrix}
\frac{1}{\sqrt{2}} & 0 & -\frac{1}{\sqrt{2}}\\[4pt]
0 & 1 & 0\\[4pt]
\frac{1}{\sqrt{2}} & 0 & \frac{1}{\sqrt{2}}
\end{bmatrix}
\begin{bmatrix}
1 & 0 & 0\\[4pt]
0 & \frac{2\sqrt{2}}{3} & \frac{1}{3}\\[4pt]
0 & -\frac{1}{3} & \frac{2\sqrt{2}}{3}
\end{bmatrix}
= \begin{bmatrix}
\frac{1}{\sqrt{2}} & \frac{\sqrt{2}}{6} & -\frac{2}{3}\\[4pt]
0 & \frac{2\sqrt{2}}{3} & \frac{1}{3}\\[4pt]
\frac{1}{\sqrt{2}} & -\frac{\sqrt{2}}{6} & \frac{2}{3}
\end{bmatrix}.
\]
\newpage

% 第5题
\noindent 5. 用Householder变换求矩阵
\[
A=\begin{bmatrix}
3 & 14 & 9 \\
6 & 43 & 3 \\
6 & 22 & 15
\end{bmatrix}
\]
的QR分解。

\noindent \textbf{答：}

\textbf{第一步：构造 $H_1$ 消去第一列后两个元素。}
取 $\bm{x}_1=[3,6,6]^{\mathsf T}$，$\|\bm{x}_1\|=\sqrt{9+36+36}=9$。
令 $\widetilde{\bm{x}}_1=[-9,0,0]^{\mathsf T}$，
\[
\bm{u}_1 = \bm{x}_1 - \widetilde{\bm{x}}_1 = \begin{bmatrix}12\\6\\6\end{bmatrix},\qquad
\|\bm{u}_1\|^2 = 144+36+36 = 216.
\]
\[
H_1 = I - 2\frac{\bm{u}_1\bm{u}_1^{\mathsf T}}{\|\bm{u}_1\|^2}
= I - \frac{1}{108}\begin{bmatrix}144 & 72 & 72\\72 & 36 & 36\\72 & 36 & 36\end{bmatrix}
= \begin{bmatrix}
-\frac13 & -\frac23 & -\frac23\\[4pt]
-\frac23 & \frac23 & -\frac13\\[4pt]
-\frac23 & -\frac13 & \frac23
\end{bmatrix}.
\]

计算 $A^{(1)} = H_1 A$：
\begin{align*}
A^{(1)}_{11} &= -\tfrac13\cdot3-\tfrac23\cdot6-\tfrac23\cdot6 = -9,\\
A^{(1)}_{21} &= -\tfrac23\cdot3+\tfrac23\cdot6-\tfrac13\cdot6 = 0,\\
A^{(1)}_{31} &= -\tfrac23\cdot3-\tfrac13\cdot6+\tfrac23\cdot6 = 0,\\
A^{(1)}_{12} &= -\tfrac13\cdot14-\tfrac23\cdot43-\tfrac23\cdot22 = -48,\\
A^{(1)}_{22} &= -\tfrac23\cdot14+\tfrac23\cdot43-\tfrac13\cdot22 = 12,\\
A^{(1)}_{32} &= -\tfrac23\cdot14-\tfrac13\cdot43+\tfrac23\cdot22 = -9,\\
A^{(1)}_{13} &= -\tfrac13\cdot9-\tfrac23\cdot3-\tfrac23\cdot15 = -15,\\
A^{(1)}_{23} &= -\tfrac23\cdot9+\tfrac23\cdot3-\tfrac13\cdot15 = -9,\\
A^{(1)}_{33} &= -\tfrac23\cdot9-\tfrac13\cdot3+\tfrac23\cdot15 = 3.
\end{align*}
\[
A^{(1)} = \begin{bmatrix}
-9 & -48 & -15\\
0 & 12 & -9\\
0 & -9 & 3
\end{bmatrix}.
\]

\textbf{第二步：构造 $H_2$ 消去 $a_{32}^{(1)}$。}
取 $\bm{x}_2=[12,-9]^{\mathsf T}$，$\|\bm{x}_2\|=\sqrt{144+81}=15$。
令 $\widetilde{\bm{x}}_2=[-15,0]^{\mathsf T}$，
\[
\bm{u}_2 = \bm{x}_2 - \widetilde{\bm{x}}_2 = \begin{bmatrix}27\\-9\end{bmatrix},\qquad
\|\bm{u}_2\|^2 = 729+81 = 810.
\]
\[
H_2' = I_2 - \frac{2}{810}\begin{bmatrix}729 & -243\\-243 & 81\end{bmatrix}
= \begin{bmatrix}-\frac45 & \frac35\\[4pt]\frac35 & \frac45\end{bmatrix}.
\]
嵌入 $3\times3$ 矩阵：
\[
H_2 = \begin{bmatrix}
1 & 0 & 0\\
0 & -\frac45 & \frac35\\[4pt]
0 & \frac35 & \frac45
\end{bmatrix}.
\]

\textbf{第三步：计算 $R$ 和 $Q$。}
\[
R = H_2 A^{(1)} = \begin{bmatrix}
-9 & -48 & -15\\
0 & -15 & 9\\
0 & 0 & -3
\end{bmatrix}.
\]

$Q^{\mathsf T}=H_2H_1$，故 $Q=H_1^{\mathsf T}H_2^{\mathsf T}=H_1H_2$（Householder矩阵对称）：
\[
Q = H_1H_2 = \begin{bmatrix}
-\frac13 & -\frac{2}{3}(-\frac45)-\frac{2}{3}(\frac35) & -\frac{2}{3}(\frac35)-\frac{2}{3}(\frac45)\\[4pt]
-\frac23 & \frac{2}{3}(-\frac45)-\frac{1}{3}(\frac35) & \frac{2}{3}(\frac35)-\frac{1}{3}(\frac45)\\[4pt]
-\frac23 & -\frac{1}{3}(-\frac45)+\frac{2}{3}(\frac35) & -\frac{1}{3}(\frac35)+\frac{2}{3}(\frac45)
\end{bmatrix}
= \begin{bmatrix}
-\frac13 & \frac{2}{15} & -\frac{14}{15}\\[4pt]
-\frac23 & -\frac{11}{15} & \frac{2}{15}\\[4pt]
-\frac23 & \frac{2}{3} & \frac{1}{3}
\end{bmatrix}.
\]
\newpage

% 第6题
\noindent 6. 用Householder变换求矩阵
\[
A=\begin{bmatrix}
0 & 4 & 1 \\
1 & 1 & 1 \\
0 & 3 & 2
\end{bmatrix}
\]
的QR分解。

\noindent \textbf{答：}

\textbf{第一步：构造 $H_1$ 消去第一列后两个元素。}
取 $\bm{x}_1=[0,1,0]^{\mathsf T}$，$\|\bm{x}_1\|=1$。
令 $\widetilde{\bm{x}}_1=[-1,0,0]^{\mathsf T}$，
\[
\bm{u}_1 = \bm{x}_1 - \widetilde{\bm{x}}_1 = \begin{bmatrix}1\\1\\0\end{bmatrix},\qquad
\|\bm{u}_1\|^2 = 2.
\]
\[
H_1 = I - 2\frac{\bm{u}_1\bm{u}_1^{\mathsf T}}{\|\bm{u}_1\|^2}
= I - \begin{bmatrix}1 & 1 & 0\\1 & 1 & 0\\0 & 0 & 0\end{bmatrix}
= \begin{bmatrix}
0 & -1 & 0\\
-1 & 0 & 0\\
0 & 0 & 1
\end{bmatrix}.
\]

计算 $A^{(1)} = H_1 A$：
\[
A^{(1)} = \begin{bmatrix}
0 & -1 & 0\\
-1 & 0 & 0\\
0 & 0 & 1
\end{bmatrix}
\begin{bmatrix}
0 & 4 & 1\\
1 & 1 & 1\\
0 & 3 & 2
\end{bmatrix}
= \begin{bmatrix}
-1 & -1 & -1\\
0 & -4 & -1\\
0 & 3 & 2
\end{bmatrix}.
\]

\textbf{第二步：构造 $H_2$ 消去 $a_{32}^{(1)}=3$。}
取 $\bm{x}_2=[-4,3]^{\mathsf T}$，$\|\bm{x}_2\|=\sqrt{16+9}=5$。
令 $\widetilde{\bm{x}}_2=[-5,0]^{\mathsf T}$，
\[
\bm{u}_2 = \begin{bmatrix}1\\3\end{bmatrix},\qquad
\|\bm{u}_2\|^2 = 10.
\]
\[
H_2' = I_2 - \frac{2}{10}\begin{bmatrix}1 & 3\\3 & 9\end{bmatrix}
= \begin{bmatrix}\frac45 & -\frac35\\[4pt]-\frac35 & -\frac45\end{bmatrix}.
\]
\[
H_2 = \begin{bmatrix}
1 & 0 & 0\\
0 & \frac45 & -\frac35\\[4pt]
0 & -\frac35 & -\frac45
\end{bmatrix}.
\]

\textbf{第三步：计算 $R$ 和 $Q$。}
\[
R = H_2 A^{(1)} = \begin{bmatrix}
-1 & -1 & -1\\
0 & \frac45(-4)-\frac35\cdot3 & \frac45(-1)-\frac35\cdot2\\[4pt]
0 & -\frac35(-4)-\frac45\cdot3 & -\frac35(-1)-\frac45\cdot2
\end{bmatrix}
= \begin{bmatrix}
-1 & -1 & -1\\
0 & -5 & -2\\
0 & 0 & -1
\end{bmatrix}.
\]

$Q = H_1H_2$：
\[
Q = \begin{bmatrix}
0 & -1 & 0\\
-1 & 0 & 0\\
0 & 0 & 1
\end{bmatrix}
\begin{bmatrix}
1 & 0 & 0\\
0 & \frac45 & -\frac35\\[4pt]
0 & -\frac35 & -\frac45
\end{bmatrix}
= \begin{bmatrix}
0 & -\frac45 & \frac35\\[4pt]
-1 & 0 & 0\\[4pt]
0 & -\frac35 & -\frac45
\end{bmatrix}.
\]

或者将符号吸收到 $R$ 中使 $R$ 对角元为正：
\[
Q = \begin{bmatrix}
0 & \frac45 & -\frac35\\[4pt]
1 & 0 & 0\\[4pt]
0 & \frac35 & \frac45
\end{bmatrix},\qquad
R = \begin{bmatrix}
1 & 1 & 1\\
0 & 5 & 2\\
0 & 0 & 1
\end{bmatrix}.
\]
\newpage

% 第7题
\noindent 7. 求矩阵
\[
A=\begin{bmatrix}
-1 & 0 & 1 & 2 \\
1 & 2 & -1 & 1 \\
2 & 2 & -2 & -1
\end{bmatrix}
\]
的满秩分解。

\noindent \textbf{答：}
满秩分解的形式为 $A = BC$，其中 $B\in\mathbb{R}^{m\times r}$ 列满秩，$C\in\mathbb{R}^{r\times n}$ 行满秩，$r=\operatorname{rank}(A)$。

\textbf{第一步：行阶梯化求秩。}
\begin{align*}
A &= \begin{bmatrix}
-1 & 0 & 1 & 2\\
1 & 2 & -1 & 1\\
2 & 2 & -2 & -1
\end{bmatrix}
\xrightarrow{R_1\leftarrow -R_1}
\begin{bmatrix}
1 & 0 & -1 & -2\\
1 & 2 & -1 & 1\\
2 & 2 & -2 & -1
\end{bmatrix}\\
&\xrightarrow{\substack{R_2\leftarrow R_2-R_1\\R_3\leftarrow R_3-2R_1}}
\begin{bmatrix}
1 & 0 & -1 & -2\\
0 & 2 & 0 & 3\\
0 & 2 & 0 & 3
\end{bmatrix}
\xrightarrow{R_3\leftarrow R_3-R_2}
\begin{bmatrix}
1 & 0 & -1 & -2\\
0 & 2 & 0 & 3\\
0 & 0 & 0 & 0
\end{bmatrix}\\
&\xrightarrow{R_2\leftarrow \frac12 R_2}
\begin{bmatrix}
1 & 0 & -1 & -2\\
0 & 1 & 0 & \frac32\\
0 & 0 & 0 & 0
\end{bmatrix}.
\end{align*}
所以 $\operatorname{rank}(A)=2$，主元列为第1列和第2列。

\textbf{第二步：构造 $B$ 和 $C$。}
取 $B$ 为 $A$ 的主元列（前两列）：
\[
B = \begin{bmatrix}
-1 & 0\\
1 & 2\\
2 & 2
\end{bmatrix}.
\]
取 $C$ 为行最简形矩阵的非零行：
\[
C = \begin{bmatrix}
1 & 0 & -1 & -2\\
0 & 1 & 0 & \frac32
\end{bmatrix}.
\]

\textbf{第三步：验证。}
\[
BC = \begin{bmatrix}
-1 & 0\\
1 & 2\\
2 & 2
\end{bmatrix}
\begin{bmatrix}
1 & 0 & -1 & -2\\
0 & 1 & 0 & \frac32
\end{bmatrix}
= \begin{bmatrix}
-1 & 0 & 1 & 2\\
1 & 2 & -1 & -2+3\\
2 & 2 & -2 & -4+3
\end{bmatrix}
= \begin{bmatrix}
-1 & 0 & 1 & 2\\
1 & 2 & -1 & 1\\
2 & 2 & -2 & -1
\end{bmatrix}
= A.
\]
\newpage

% 第8题
\noindent 8. 求矩阵
\[
A=\begin{bmatrix}
1 & 2 & 3 & 0 \\
0 & 2 & 1 & -1 \\
1 & 0 & 2 & 1
\end{bmatrix}
\]
的满秩分解。

\noindent \textbf{答：}

\textbf{第一步：行阶梯化求秩。}
\begin{align*}
A &= \begin{bmatrix}
1 & 2 & 3 & 0\\
0 & 2 & 1 & -1\\
1 & 0 & 2 & 1
\end{bmatrix}
\xrightarrow{R_3\leftarrow R_3-R_1}
\begin{bmatrix}
1 & 2 & 3 & 0\\
0 & 2 & 1 & -1\\
0 & -2 & -1 & 1
\end{bmatrix}\\
&\xrightarrow{R_3\leftarrow R_3+R_2}
\begin{bmatrix}
1 & 2 & 3 & 0\\
0 & 2 & 1 & -1\\
0 & 0 & 0 & 0
\end{bmatrix}
\xrightarrow{R_2\leftarrow \frac12 R_2}
\begin{bmatrix}
1 & 2 & 3 & 0\\
0 & 1 & \frac12 & -\frac12\\
0 & 0 & 0 & 0
\end{bmatrix}\\
&\xrightarrow{R_1\leftarrow R_1-2R_2}
\begin{bmatrix}
1 & 0 & 2 & 1\\
0 & 1 & \frac12 & -\frac12\\
0 & 0 & 0 & 0
\end{bmatrix}.
\end{align*}
所以 $\operatorname{rank}(A)=2$，主元列为第1列和第2列。

\textbf{第二步：构造 $B$ 和 $C$。}
\[
B = \begin{bmatrix}
1 & 2\\
0 & 2\\
1 & 0
\end{bmatrix},\qquad
C = \begin{bmatrix}
1 & 0 & 2 & 1\\
0 & 1 & \frac12 & -\frac12
\end{bmatrix}.
\]

\textbf{第三步：验证。}
\[
BC = \begin{bmatrix}
1 & 2\\
0 & 2\\
1 & 0
\end{bmatrix}
\begin{bmatrix}
1 & 0 & 2 & 1\\
0 & 1 & \frac12 & -\frac12
\end{bmatrix}
= \begin{bmatrix}
1 & 2 & 2+1 & 1-1\\
0 & 2 & 0+1 & 0-1\\
1 & 0 & 2+0 & 1+0
\end{bmatrix}
= \begin{bmatrix}
1 & 2 & 3 & 0\\
0 & 2 & 1 & -1\\
1 & 0 & 2 & 1
\end{bmatrix}
= A.
\]
\newpage

% 第9题
\noindent 9. 求矩阵
\[
A=\begin{bmatrix}
1 & 0 & 1 \\
0 & 1 & 1 \\
0 & 0 & 0
\end{bmatrix}
\]
的奇异值分解。

\noindent \textbf{答：}
SVD 的形式为 $A = U\Sigma V^{\mathsf T}$。

\textbf{第一步：计算 $A^{\mathsf T}A$ 的特征值和特征向量。}
\[
A^{\mathsf T}A = \begin{bmatrix}
1 & 0 & 0\\
0 & 1 & 0\\
1 & 1 & 0
\end{bmatrix}
\begin{bmatrix}
1 & 0 & 1\\
0 & 1 & 1\\
0 & 0 & 0
\end{bmatrix}
= \begin{bmatrix}
1 & 0 & 1\\
0 & 1 & 1\\
1 & 1 & 2
\end{bmatrix}.
\]

求特征值：$|\lambda I - A^{\mathsf T}A| = 0$。
\begin{align*}
\begin{vmatrix}
\lambda-1 & 0 & -1\\
0 & \lambda-1 & -1\\
-1 & -1 & \lambda-2
\end{vmatrix}
&= (\lambda-1)[(\lambda-1)(\lambda-2)-1] - [0-(\lambda-1)]\\
&= (\lambda-1)(\lambda^2-3\lambda+1) - (\lambda-1)\\
&= (\lambda-1)(\lambda^2-3\lambda)\\
&= \lambda(\lambda-1)(\lambda-3) = 0.
\end{align*}
特征值为 $\lambda_1=3,\;\lambda_2=1,\;\lambda_3=0$。故奇异值为 $\sigma_1=\sqrt{3},\;\sigma_2=1,\;\sigma_3=0$。

\textbf{第二步：求 $V$（$A^{\mathsf T}A$ 的特征向量）。}

\textit{对 $\lambda_1=3$：}解 $(3I-A^{\mathsf T}A)\bm{v}_1=\bm{0}$。
\[
\begin{bmatrix}
2 & 0 & -1\\
0 & 2 & -1\\
-1 & -1 & 1
\end{bmatrix}\begin{bmatrix}x\\y\\z\end{bmatrix}=\bm{0}
\;\Rightarrow\;
\begin{cases}2x=z\\2y=z\\x+y=z\end{cases}
\;\Rightarrow\; x=y,\;z=2x.
\]
取 $\bm{v}_1=(1,1,2)^{\mathsf T}$，单位化得 $\bm{v}_1=\frac{1}{\sqrt{6}}(1,1,2)^{\mathsf T}$。

\textit{对 $\lambda_2=1$：}
\[
\begin{bmatrix}
0 & 0 & -1\\
0 & 0 & -1\\
-1 & -1 & -1
\end{bmatrix}\begin{bmatrix}x\\y\\z\end{bmatrix}=\bm{0}
\;\Rightarrow\; z=0,\;x+y=0.
\]
取 $\bm{v}_2=(1,-1,0)^{\mathsf T}$，单位化得 $\bm{v}_2=\frac{1}{\sqrt{2}}(1,-1,0)^{\mathsf T}$。

\textit{对 $\lambda_3=0$：}
\[
\begin{bmatrix}
-1 & 0 & -1\\
0 & -1 & -1\\
-1 & -1 & -2
\end{bmatrix}\begin{bmatrix}x\\y\\z\end{bmatrix}=\bm{0}
\;\Rightarrow\; x=-z,\;y=-z.
\]
取 $\bm{v}_3=(1,1,-1)^{\mathsf T}$，单位化得 $\bm{v}_3=\frac{1}{\sqrt{3}}(1,1,-1)^{\mathsf T}$。
\[
V = \begin{bmatrix}
\frac{1}{\sqrt{6}} & \frac{1}{\sqrt{2}} & \frac{1}{\sqrt{3}}\\[4pt]
\frac{1}{\sqrt{6}} & -\frac{1}{\sqrt{2}} & \frac{1}{\sqrt{3}}\\[4pt]
\frac{2}{\sqrt{6}} & 0 & -\frac{1}{\sqrt{3}}
\end{bmatrix},\qquad
\Sigma = \begin{bmatrix}
\sqrt{3} & 0 & 0\\
0 & 1 & 0\\
0 & 0 & 0
\end{bmatrix}.
\]

\textbf{第三步：求 $U$。}
对于非零奇异值：$\bm{u}_i = \frac{1}{\sigma_i}A\bm{v}_i$。
\begin{align*}
\bm{u}_1 &= \frac{1}{\sqrt{3}}\begin{bmatrix}1 & 0 & 1\\0 & 1 & 1\\0 & 0 & 0\end{bmatrix}
\frac{1}{\sqrt{6}}\begin{bmatrix}1\\1\\2\end{bmatrix}
= \frac{1}{\sqrt{18}}\begin{bmatrix}1+2\\1+2\\0\end{bmatrix}
= \frac{1}{\sqrt{2}}\begin{bmatrix}1\\1\\0\end{bmatrix},\\[4pt]
\bm{u}_2 &= \frac{1}{1}\begin{bmatrix}1 & 0 & 1\\0 & 1 & 1\\0 & 0 & 0\end{bmatrix}
\frac{1}{\sqrt{2}}\begin{bmatrix}1\\-1\\0\end{bmatrix}
= \frac{1}{\sqrt{2}}\begin{bmatrix}1\\-1\\0\end{bmatrix}.
\end{align*}

$\bm{u}_3$ 需与 $\bm{u}_1,\bm{u}_2$ 构成标准正交基，可取 $\bm{u}_3=(0,0,1)^{\mathsf T}$。
\[
U = \begin{bmatrix}
\frac{1}{\sqrt{2}} & \frac{1}{\sqrt{2}} & 0\\[4pt]
\frac{1}{\sqrt{2}} & -\frac{1}{\sqrt{2}} & 0\\[4pt]
0 & 0 & 1
\end{bmatrix}.
\]

\textbf{第四步：写出最终 SVD。}
\[
A = U\Sigma V^{\mathsf T} = \begin{bmatrix}
\frac{1}{\sqrt{2}} & \frac{1}{\sqrt{2}} & 0\\[4pt]
\frac{1}{\sqrt{2}} & -\frac{1}{\sqrt{2}} & 0\\[4pt]
0 & 0 & 1
\end{bmatrix}
\begin{bmatrix}
\sqrt{3} & 0 & 0\\
0 & 1 & 0\\
0 & 0 & 0
\end{bmatrix}
\begin{bmatrix}
\frac{1}{\sqrt{6}} & \frac{1}{\sqrt{6}} & \frac{2}{\sqrt{6}}\\[4pt]
\frac{1}{\sqrt{2}} & -\frac{1}{\sqrt{2}} & 0\\[4pt]
\frac{1}{\sqrt{3}} & \frac{1}{\sqrt{3}} & -\frac{1}{\sqrt{3}}
\end{bmatrix}.
\]

\end{document}